\section*{Introduction}
\label{s:introduction}
Bone is a living composite that continually reorganizes its microstructure to balance mechanical function with metabolic cost. Four mechanobiologically distinct processes drive this adaptation: (i) the elastic response of the solid matrix to loads, (ii) the evolution of mass (apparent density/porosity), (iii) the reorientation of the trabecular/collagen network (fabric anisotropy), and (iv) the production and spatial transport of cell-mediated biochemical signals that couple local mechanics to osteoblast/osteoclast activity. Robust numerical prediction of morphology and stiffness therefore requires a model that co-evolves these \emph{four} fields---displacement $u$, density $\rho$, fabric tensor $\mathbf{A}$, and a diffusive mechanotransductive stimulus $S$---within a single, tightly coupled formulation.

Mechanobiology supports this coupling. First, the mechanical field is not merely a passive input: Turner’s “three rules” assert that the magnitude, rate (frequency), and rest intervals of \emph{dynamic} loading govern the osteogenic response \cite{Turner1998_ThreeRulesBone}. Spatial \emph{strain gradients}---proxies for lacuno--canalicular (LCN) fluid flow---localize bone formation more reliably than strain magnitude alone, as shown in turkey radius models \cite{Gross1997_StrainGradientsCorrelate} and in exercise-induced periosteal activation in rooster models \cite{Zernicke1997_StrainGradientsCorrelate}. These gradients act through osteocytes embedded in the LCN, where fluid shear induces Ca$^{2+}$/Wnt signaling; acute loading suppresses sclerostin (Sost), while unloading elevates it, shifting the local formation/resorption balance \cite{Robling2008_MechanicalStimulationBone}. While mechanosensitive channels such as PIEZO1 contribute to skeletal homeostasis, their cell-type-specific roles in load adaptation are nuanced and do not obviate the need to resolve the full biophysical cascade from matrix strains to fluid shear to osteocytic signaling \cite{Wang2020_MechanicalSensingProtein}. Together, these findings motivate an explicit stimulus field $S$ that is \emph{generated} by mechanical cues and \emph{transported} over osteocytic/vascular pathways (e.g., via diffusion--decay dynamics), rather than a purely local, algebraic “mechanical index.”

Second, accurate mechanics--morphology feedback requires density $\rho$ and anisotropy $\mathbf{A}$ to co-evolve and feed back into the elastic moduli. Empirically, bone volume fraction (BV/TV) explains most variance in cancellous stiffness, and fabric anisotropy accounts for nearly all of the remainder; together they reach $\sim$95--98\% explanatory power across anatomical sites, with negligible contribution from other morphometric covariates \cite{Maquer2015_BoneVolumeFraction, Maquer2017_NotOnlyStiffness}. Theoretical and constitutive links between $\{\rho,\mathbf{A}\}$ and orthotropic elastic constants are well established---Cowin’s fabric--elasticity relations (1985, 1990) and the Zysset--Curnier orthotropic model (1995) connect BV/TV and second-order fabric tensors to the apparent stiffness tensor (including $E$, $G$, and $\nu$) \cite{Cowin1985_RelationshipBetweenElasticity, Turner1990_FabricDependenceOrthotropic, Zysset1995_AlternativeModelAnisotropic}. In practice, fabric can be quantified directly from imaging using gray-level structure-tensor or mean-intercept-length methods and validated on HR-pQCT/$\mu$CT, enabling evolution laws for $\mathbf{A}$ aligned with principal stress/flow directions \cite{Kuo2007_QuantifyingAnisotropyTrabecular, Ali2017_Mu, Rietbergen2015_NovelApproachEstimate, Dutra2016_QuantifyingTrabecularBone}. As $\rho$ and $\mathbf{A}$ change through remodeling, the spatial stiffness field---and thus the displacement field $u$ and its gradients---must be updated consistently to avoid spurious transients and mesh-dependent artifacts; morphology--elasticity models with full fabric information achieve the highest predictive accuracy in such coupled simulations \cite{Schneider2012_MorphologyElasticityRelationships, Panyasantis2016_EffectiveElasticProperties}.

Third, the remodeling stimulus that drives BMU-scale biology is not localized to a voxel or finite element. Experiments and theory both indicate spatial nonlocality through osteocytic/vascular networks; modern formulations therefore promote the stimulus to a \emph{field} obeying a diffusion--reaction equation, which systematically removes unphysical checkerboarding and captures healing or damage--remodeling interactions \cite{Giorgio2019_MechanicallyDrivenBiological, Addessi2024_BoneRemodelingApproach, Kobler2024_GradientEnhancedBone}. Recent work also explores multi-stimulus energetics (e.g., shear vs.\ normal components) in continuum settings, reinforcing the need to track multiple evolving material parameters beyond a single density variable \cite{Branecka2024_BoneRemodelingModel}.

From a numerical standpoint, loosely coupled or explicit updates are in general more fragile than monolithic schemes: density-only algorithms are prone to mesh-dependent checkerboards and near-field oscillations near load singularities \cite{Jacobs1995_NumericalInstabilitiesBone, Kobler2024_GradientEnhancedBone}. We therefore retain a partitioned block Gauss--Seidel fixed point and wrap it in Anderson acceleration with scaled Tikhonov regularization, step limiting, and adaptive history restarts; see the solver description in \S\ref{sec:fp-aa}. This combination substantially mitigates splitting artifacts and yields robust, quasi mesh-independent convergence in our experiments \cite{Both2019_AndersonAcceleratedFixed, Walker2011_AndersonAcceleration, Fang2009_TwoClassesMultisecant} while remaining minimally invasive on top of standard FE components.

